\documentclass[UTF8,a4paper,11pt]{ctexart}
\usepackage[a4paper,margin=2.1cm]{geometry}
\usepackage{amsmath,booktabs,longtable,tabularx,array,tikz,xcolor,hyperref,enumitem,fancyhdr,tcolorbox}
\usetikzlibrary{arrows.meta,positioning,fit,shapes.geometric,calc}
\definecolor{navy}{HTML}{16324F}\definecolor{blue}{HTML}{2F6690}\definecolor{orange}{HTML}{E07A5F}\definecolor{lightblue}{HTML}{EAF3F8}\definecolor{lightorange}{HTML}{FCEEEA}
\hypersetup{colorlinks=true,linkcolor=blue,urlcolor=blue,citecolor=blue}
\pagestyle{fancy}\fancyhf{}\lhead{具身智能技术栈}\rhead{VLA / WAM / Code-as-Policy}\cfoot{\thepage}
\setlist{nosep,leftmargin=2em}
\newcommand{\source}[1]{\href{#1}{[来源]}}
\newtcolorbox{insight}{colback=lightblue,colframe=blue,boxrule=.5pt,arc=2pt}
\newtcolorbox{warning}{colback=lightorange,colframe=orange,boxrule=.5pt,arc=2pt}
\title{具身智能当前技术栈\\[-2mm]\large VLA、WAM 与 Code-as-Policy 的差异、接口和进展}
\author{技术补充讲义（资料冻结：2026 年 9 月 12 日）}\date{}
\begin{document}\maketitle\tableofcontents\newpage

\section{先看系统：模型只是中间一层}
一个可部署的机器人系统至少包含传感器、状态估计、基础模型、动作转换、低层控制和安全监控。三条路线主要差在“基础模型输出什么”，而不在是否需要相机、标定、动力学和急停。

\begin{figure}[ht]
\centering
\begin{tikzpicture}[>=Latex,node distance=8mm and 8mm,
box/.style={draw=navy,rounded corners,fill=lightblue,align=center,minimum width=2.65cm,minimum height=10mm},
low/.style={draw=navy,rounded corners,fill=white,align=center,minimum width=2.65cm,minimum height=10mm},
arr/.style={->,thick,draw=navy}]
\node[box](sense){相机 / 深度 / 力\\\\本体状态};
\node[box,right=of sense](rep){视觉与语言\\\\表征};
\node[box,right=of rep](fm){VLA / WAM / Code-as-Policy};
\node[box,right=of fm](trans){动作解码\\位姿、关节、程序};
\node[low,right=of trans](ctrl){IK、轨迹规划\\伺服与安全};
\node[low,below=of ctrl](robot){机器人身体\\执行与反馈};
\draw[arr](sense)--(rep)--(fm)--(trans)--(ctrl)--(robot);
\draw[arr](robot.south west) |- node[pos=.25,left]{观测闭环}(sense.south);
\node[draw=orange,dashed,fit=(rep)(fm)(trans),inner sep=5mm,label={[orange]above:学习模型层}] {};
\end{tikzpicture}
\caption{部署系统的共同骨架。VLA、WAM 和 Code-as-Policy 不能单独替代 IK、伺服控制、标定或安全层。}
\end{figure}

\section{三条路线的核心差异}
\begin{longtable}{p{2.8cm}p{3.4cm}p{3.4cm}p{3.4cm}}
\toprule
维度 & VLA & WAM & Code-as-Policy\\\midrule
模型输出 & 动作 token 或连续 action chunk & 未来视频 latent 与动作 chunk & 可执行程序、API 调用或几何约束\\
主要监督 & 机器人示范、语言、视觉预训练 & 视频先验加机器人/人类动作 & 语言、技能库、执行轨迹和反馈\\
推理方式 & 一次或滚动动作块解码 & 联合去噪，预测“将发生什么”和“怎么做” & 生成、执行、观察、修复的循环\\
优势 & 低层动作连贯、端到端程度高 & 具备物理未来建模和数据扩展能力 & 可解释、可组合、易适配 API\\
瓶颈 & 本体差异、数据规模、长时记忆 & 视频逼真不等于动作可执行 & API 安全、延迟、接触控制和错误累积\\
代表系统 & RT-2、$\pi_0$、Xiaomi-Robotics-0/1 & OpenWAM-$\alpha$、Xiaomi-Robotics-U0 & CaP、CaP-X、ASPIRE、PhysCaP\\
\bottomrule
\end{longtable}

\section{VLA 技术栈与进展}
\begin{figure}[ht]
\centering
\begin{tikzpicture}[>=Latex,node distance=8mm and 11mm,
box/.style={draw=navy,rounded corners,fill=lightblue,align=center,minimum width=2.9cm,minimum height=10mm},arr/.style={->,thick,draw=navy}]
\node[box](web){互联网图文\\\\语言知识};\node[box,right=of web](vlm){视觉语言骨干};\node[box,right=of vlm](expert){动作专家token / flow};\node[box,right=of expert](chunk){动作块};\node[box,right=of chunk](body){本体适配+ 伺服};
\draw[arr](web)--(vlm)--(expert)--(chunk)--(body);
\node[below=of expert,align=center]{RT-2：动作 token\\$\pi_0$：连续 flow matching};
\end{tikzpicture}
\caption{VLA 的典型数据流。RT-2 重点在语言知识迁移，$\pi_0$ 重点在连续动作生成。}\source{https://arxiv.org/abs/2307.15818}\;\source{https://www.pi.website/blog/pi0}
\end{figure}

VLA 的关键工程选择包括：动作是离散 token 还是连续向量；预测单步还是 chunk；是否使用历史帧；如何统一不同机器人自由度；以及高层语言条件是否与低层动作共享骨干。$\pi_{0.5}$、$\pi^*_{0.6}$ 和 $\pi_{0.7}$ 代表从开放环境泛化、在线强化学习到可操纵条件的连续推进。小米 Xiaomi-Robotics-0/1 则强调实时执行和大规模真实轨迹。

\section{WAM 技术栈与进展}
WAM 把世界模型和动作模型的训练目标耦合起来。OpenWAM-$\alpha$ 使用视频 DiT 与专用 ActionDiT，通过联合注意力和同步去噪，使动作影响的未来帧与动作本身保持一致。\source{https://openwam-official.github.io/}

\[
 \mathcal{L}=\mathcal{L}_{\mathrm{video}}+\lambda\mathcal{L}_{\mathrm{action}}+\gamma\mathcal{L}_{\mathrm{coupling}},
 \qquad
 (\hat{x}_{1:T},\hat{a}_{1:K})=f_\theta(x_0,L,\epsilon_x,\epsilon_a).
\]

它的价值不是只生成好看的视频，而是让视频预测成为动作选择的可检验中间表征。风险在于世界模型的视觉准确度、动作可执行性和真实动力学之间仍可能存在偏差。Xiaomi-Robotics-U0 应放在同一“世界模型和数据引擎”分支，而不是与直接控制 VLA 混为一谈。

\section{Code-as-Policy 技术栈与进展}
\begin{figure}[ht]
\centering
\begin{tikzpicture}[>=Latex,node distance=8mm and 12mm,
box/.style={draw=navy,rounded corners,fill=lightblue,align=center,minimum width=2.8cm,minimum height=10mm},arr/.style={->,thick,draw=navy}]
\node[box](prompt){目标、示例\\和约束};\node[box,right=of prompt](gen){LLM 生成\\\\程序 / 技能};\node[box,right=of gen](check){静态检查\\\\类型、碰撞、权限};\node[box,right=of check](exec){API 执行IK / 轨迹 / 伺服};\node[box,below=of exec](fb){视觉差分、状态、\\\\成功/失败反馈};
\draw[arr](prompt)--(gen)--(check)--(exec);\draw[arr](exec)--(fb);\draw[arr](fb.west) -| (gen.south);
\end{tikzpicture}
\caption{Code-as-Policy 的工程闭环。CaP-X、ASPIRE 和 PhysCaP 的进展主要发生在反馈、技能复用和探索，而不是单纯扩大提示词。}\source{https://code-as-policies.github.io/}\;\source{https://capgym.github.io/}
\end{figure}

CaP 的接口设计决定系统上限：API 越高层，生成越稳定但表达能力越有限；API 越接近关节或力矩，表达能力越强但安全验证和实时性越困难。Astra 的 Robocurve 实验使用末端位姿 API，再由 IK 转为关节角，正是这种折中。

\section{GPT-6 Astra 应放在哪里}
GPT-6 Astra 的官方定位是通用推理、代码和 computer use 模型，并未公开声明为专门 VLA。\source{https://openai.com/index/gpt-6-astra/} Robocurve 的 2026-09-04 测试表明，配合 Inspect Robots harness，它可以在双臂 YAM 上完成高层视觉闭环控制：方块入碗 19/20，拼图插槽 2/20。\source{https://openai.robocurve.org/gpt-6-astra/}

因此其技术栈位置是
\[
 \text{通用多模态模型}
 +\text{机器人 API}
 +\text{视觉/状态反馈}
 +\text{IK 与安全控制}
 \quad\Rightarrow\quad
 \text{Code-as-Policy Agent}.
\]
这类结果很有启发性：通用模型的空间推理和代码能力可以迁移到物理控制；但有限任务、不同机械臂、人工评分和精细接触失败都说明它还不能替代专门 VLA。X 上的 physical ICL 演示和小红书视频可作为定性案例，不能替代标准化实验。\source{https://www.sotwe.com/npew?lang=en}\;\source{https://www.yicai.com/news/103359498.html}

\section{模型选型与评测清单}
\begin{enumerate}
\item 需要低延迟连续动作和多机器人控制：优先比较 VLA，报告 action chunk、控制频率和本体适配。
\item 需要预测后果、生成数据或 OOD 场景：比较 WAM，检查视频质量与真实动作成功率是否一致。
\item 需要快速组合新技能、调用现有机器人软件：比较 Code-as-Policy，限制 API 权限并记录每次反馈修复。
\item 任何路线都要报告成功率、阶段得分、失败类型、耗时、成本、硬件、是否微调和测试环境；社交媒体视频只进入证据等级为“演示”的附录。
\end{enumerate}

\bibliographystyle{unsrt}\bibliography{references}
\end{document}







